\scanpage{26}
\deleted{Si $A$ est l'algèbre de von Neumann produit $\prod$ des $A_i$ $(i\in I)$ et d'une algèbre de von Neumann}

Liebe Paula\unsure{!|,}

\deleted{$B$ \inserted{sans projecteur}}

Wieder ein Breifn,

\deleted{minimal. $T$ induit sur $A_i$ la fonction $\lambda_i T_i$, où $\lambda_i>0$ et $T_i$ est la trace usuelle.}

\deleted{Mais les $A_i$ peuvent se plonger dans une}

\deleted{\unsure{Une}} Une algèbre de von Neumann \deleted{quelconque} $A$ \inserted{peut être considérée comme} \deleted{peut se plonger dans une} algèbre de von Neumann \inserted{d'une algèbre de von Neumann} $B$ non atomique, \deleted{\unsure{par exemple}}, et si $A$ est muni d'une trace semi-finie fidèle $T$, on peut choisir un $B$ muni d'une trace semi-finie fidèle qui induit \deleted{\unsure{exactement}} sur $A$ la trace donnée.

Par exemple, on prend l'algèbre des applications, \deleted{\unsure{bornées}} faiblement mesurables et bornées de $[0,1]$ dans $A\subset L(\mathcal H)$, \deleted{\unsure{munie de}} (sur, en réalité sur $\mathcal H\otimes\mathcal H'$) avec la trace
\[
\overline T(f)=\int_0^1 T(f(s))\,ds.
\]
Bien entendu, si $B$ est une algèbre de von Neumann munie d'une trace $T$, et $A$ une sous-algèbre de von Neumann, alors $A\cap\mathfrak A_2$ \deleted{d\unsure{ans}} dans l'adhérence \inserted{alors} l'idéal de définition de la trace sur $A$, $\mathfrak H_A$ est la \unsure{même} finie. Considérons $A$ une rétraction de $A$ sur $B$, on voit \deleted{que dans les \unsure{formes} positives} \inserted{dans le plongement en question, sur les traces} finies, \unsure{tout|toute} la \unsure{théorème|théorie} \unsure{encore|connue} sur $A$ est analogue.
\mspage{27}
